\documentclass{article}

\usepackage{fullpage}
\usepackage{url}
\usepackage{graphicx}
\usepackage{amsmath}
\usepackage{physics}
\usepackage{mathtools}
\usepackage{bbold}

\DeclareMathOperator{\erfi}{erfi}
\newcommand{\R}{\ensuremath{\mathbb{R}}}
\newcommand{\C}{\ensuremath{\mathbb{C}}}

\title{Statistical Mechanics I Excercise 6}
\author{Idan Stark}
\date{\today}

\begin{document}
\maketitle

\section{Landau Theory for Complex Valued Fields}
Let us consider the following Landau free energy:
\begin{equation}
    \beta\mathcal{F}\left[\psi(x),\phi(x)\right] = \int d^dx\left(
        \frac{K}{2}\left|\nabla\psi\right|^2 + t\left|\psi\right|^2 + u\left|\psi\right|^4 + v\left|\psi\right|^6 + \frac{\phi^2}{2\sigma} - \gamma\sigma\left|\psi\right|^2
    \right)
\end{equation}

\subsection{Effective theory}
Let us look at the partition function for that energy:
\begin{equation}
\begin{gathered}
    Z = \int D[\psi(x), \phi(x)]\exp[-\beta\mathcal{F}] = \\ =
    \int D[\psi(x), \phi(x)] \exp[-\int d^dx\left(
        \frac{K}{2}\left|\nabla\psi\right|^2 + t\left|\psi\right|^2 + u\left|\psi\right|^4 + v\left|\psi\right|^6 + \frac{\phi^2}{2\sigma} - \gamma\phi\left|\psi\right|^2
    \right)] = \\ =
    \int D[\psi(x)] 
    \Biggl(
    \exp[-\int d^dx\left(
        \frac{K}{2}\left|\nabla\psi\right|^2 + t\left|\psi\right|^2 + u\left|\psi\right|^4 + v\left|\psi\right|^6
    \right)] \times \\
    \int D[\phi(x)]
    \exp[-\int d^dx\left(
        \frac{\phi^2}{2\sigma} - \gamma\phi\left|\psi\right|^2
    \right)]
    \Biggr)
\end{gathered}
\end{equation}
The inner integral is the path intergral equivilent of a gaussian, and can be calculated by writing it explicitly and pulling the product out of the integral:
\begin{equation}
\begin{gathered}
    \int D[\phi(x)]
    \exp[-\int d^dx\left(
        \frac{\phi^2}{2\sigma} - \gamma\phi\left|\psi\right|^2
    \right)] = \\ =
    \int \prod_x d\phi(x) \exp[-\left(
        \frac{\phi(x)^2}{2\sigma} - \gamma\phi(x)\left|\psi\right|^2
    \right)] = \\ =
     \prod_x \int d\phi(x) \exp[-\left(
        \frac{\phi(x)^2}{2\sigma} - \gamma\phi(x)\left|\psi\right|^2
    \right)] = \\ =
    \prod_x e^{\frac{\gamma^2\sigma}{2}\left|\psi\right|^4}\int d\phi(x) \exp[-\left(
        \frac{\phi(x)^2}{2\sigma} - \gamma\phi(x)\left|\psi\right|^2 + \frac{\gamma^2\sigma}{2}\left|\psi\right|^4
    \right)] = \\ =
    \prod_x e^{\frac{\gamma^2\sigma}{2}\left|\psi\right|^4}\int d\phi(x) \exp[-\frac{\left(\phi(x) - \gamma\sigma\left|\psi\right|^2\right)^2}{2\sigma}] = \\ =
    e^{\int d^dx \frac{\gamma^2\sigma}{2}\left|\psi\right|^4}
    \prod_x \left(\sigma \pi\right) = \\ =
    \exp[\int d^dx \left(
        \frac{\gamma^2\sigma}{2}\left|\psi\right|^4 + \ln(\sigma\pi)
    \right)]
\end{gathered}
\end{equation}
This means the integral over $\phi(x)$ added to the effective $u$ in the theory and a constant energy that can be ignored. The effective free energy is then:
\begin{equation}
    \beta \mathcal{F}[\psi(x)] = \int d^dx\left(
        \frac{K}{2}\left|\nabla\psi\right|^2 + t\left|\psi\right|^2 + \left(u - \frac{\gamma^2\sigma}{2}\right)\left|\psi\right|^4 + v\left|\psi\right|^6
    \right)
\end{equation}
\subsection{Phase Diagram}
Let us now use the new energy to find the phase diagram of the effective theory around the minima of the free energy. At that area, using the saddle point approximation, the fluctuations in $\psi$ are much smaller then the field itself $\psi$. This means we can neglect the first term and focus on the other three. The three remaining terms give the same Landau free energy we saw in assignment 4, which gave a tricrticial point due to the $\psi^6$ term (there is was $m^6$).
% TODO phase diagram

\section{Non-Homogeneous Field Model}
Some d dimensions system is described by an order parameter that has a special
axis, $x_{||}$, namely $m(x_{||}, \vec{x}_{\perp})$. Its Landau-Ginzburg free energy is given by:
\begin{equation}
    \beta \mathcal{F}[m(x)] = \int dx_{||} \int dx^{d-1}_{\perp} \left[
        \frac{t}{2}m^2 + \frac{K}{2}\left(\partial_{x_{||}} m\right)^2  + \frac{L}{2}\left(\nabla_{\perp}^2 m\right)^2 - hm
    \right]
\end{equation}
where $\nabla_{\perp}^2 = \sum_{x_i \ne x_{||}}\partial_{x_i}^2$.
\subsection{Symmetries}
The first symmtery we notice wehn $h = 0$ is the $\mathbb{Z}_2$ symmetry $m \rightarrow -m$. The second group of symmetries is the Orthogonal group $O(d - 1)$ of rotation and mirroring symmetries for the $d - 1$ dimensions $x_i \ne x_{||}$. Lastly, there is a mirroring $\mathbb{Z}_2$ symmetry for the mirroring along $x_{||}$.
\subsection{Fourier Space Free energy}
Let us now transform the energy to Fourier space. Writing:
\begin{equation}
    m = \int \frac{d^dk}{(2\pi)^d} \tilde{m}(\vec{k})e^{i\vec{k}\cdot \vec{x}}
\end{equation}
The free energy functional becomes:
\begin{equation}
\begin{gathered}
    \beta \mathcal{F}[\tilde{m}(\vec{k})] = \int dx_{||} \int dx^{d-1}_{\perp} \left[
        \frac{t}{2}m^2 + \frac{K}{2}\left(\partial_{x_{||}} m\right)^2  + \frac{L}{2}\left(\nabla_{\perp}^2 m\right)^2 - hm
    \right] = \\ = \int dx_{||} \int dx^{d-1}_{\perp} \int \frac{d^dk}{(2\pi)^d} \Biggr[
        \int \frac{d^dk'}{(2\pi)^d} 
        \Biggr(
            \frac{t}{2}\tilde{m}(\vec{k})\tilde{m}(\vec{k}')e^{i(\vec{k} + \vec{k}')\cdot\vec{x}} + \\  + \frac{K}{2}\left(\partial_{x_{||}}\tilde{m}(\vec{k})e^{i\vec{k}\cdot\vec{x}}\right)\left(\partial_{x_{||}}\tilde{m}(\vec{k}')e^{i\vec{k}'\cdot\vec{x}}\right)+ \frac{L}{2}\left(\nabla_{\perp}^2 \tilde{m}(\vec{k})\tilde{m}(\vec{k})e^{i\vec{k}\cdot\vec{x}}\right)\left(\nabla_{\perp}^2 \tilde{m}(\vec{k})\tilde{m}(\vec{k}')e^{i\vec{k}'\cdot\vec{x}}\right)
        \Biggr) - h\tilde{m}(\vec{k})e^{i\vec{k}\cdot \vec{x}}
    \Biggr]
\end{gathered}
\end{equation}
Each of the terms where there are two $\tilde{m}$, evaluating the space integral over the exponents gives a delta function $(2\pi)^d \delta(\vec{k} + \vec{k}')$, while last term that has only one exponent gives $(2\pi)^d\delta(\vec{k})$. So, using these delta functions to integrate over the relevant $k$ integrals gives:
\begin{equation}
    \beta \mathcal{F}[\tilde{m}(k_{||}, \vec{k}_{\perp})] = \int \frac{d^dk}{(2\pi)^d}
    \tilde{m}(k_{||}, \vec{k}_{\perp})\tilde{m}(-k_{||}, -\vec{k}_{\perp})
    \left(
        \frac{t}{2} + \frac{K}{2}k_{||}^2 + \frac{L}{2}\vec{k}_{\perp}^4
    \right) - h\tilde{m}(0, \vec{0})
\end{equation}
\subsection{RG Recursion relations}
Let us now find the recursion relations for the constants in the theory. For this, let us rescale $k_{||}' = bk_{||}$ and $\vec{k}_{\perp}' = c\vec{k}_{\perp}$ and renormalize the field $\tilde{m}' = m/z$. This gives:
\begin{equation}
    \beta \mathcal{F}[\tilde{m}(k_{||}', \vec{k}_{\perp}')] = \int \frac{d^dk'}{(2\pi)^d}\frac{z^2}{bc^{d-1}}
    \tilde{m}(k_{||}', \vec{k}_{\perp}')\tilde{m}(-k_{||}', -\vec{k}_{\perp}')
    \left(
        \frac{t}{2} + \frac{K}{2b^2}k_{||}'^2 + \frac{L}{2c^4}\vec{k}_{\perp}'^4
    \right) - h\tilde{m}'(0, \vec{0})
\end{equation}
This gives the recursion relations:
\begin{equation}
\begin{gathered}
    t'/t = z^2b^{-1}c^{1-d}
    \\
    K'/K = z^2b^{-3}c^{1-d}
    \\
    L'/L = z^2b^{-1}c^{-3-d}
\end{gathered}
\end{equation}
\subsection{Requiring unchanged prefactors}
Next, let us require that $K' = K$ and $L' = L$. This gives:
\begin{equation}
\begin{cases}
    z^2b^{-3}c^{1-d} = 1
    \\
    z^2b^{-1}c^{-3-d} = 1
\end{cases}
\end{equation}
Solving the equations for $c(b)$ and $z(b)$ gives:
\begin{equation}
    c(b) = b^{1/2}
    \qquad
    z = b^{(d + 5)/4}
\end{equation}
This means that:
\begin{equation}
    t' = b^2t
\end{equation}
\subsection{Critical Exponents}
Let us finally find the critical exponent $\nu$, for the correlation length:
\begin{equation}
    \xi(t) \sim t^{-\nu}
\end{equation}
At ever RG space, the correlation length scales according to the constants $b$ and $c$:
\begin{equation}
    \xi_{||}'(t') = b^{-1}\xi_{||}(t)
    \qquad
    \xi_{\perp}'(t') = c^{-1}\xi(t) = b^{-1/2}\xi_{\perp}(t)
\end{equation}
So:
\begin{equation}
    t'^{-\nu_{||}} = b^{-2\nu_{||}}t^{-\nu_{||}} \sim \xi_{||}'(t') = b^{-1}\xi_{||}(t) \sim b^{-1}t^{-\nu_{||}} \rightarrow \nu_{||} = \frac{1}{2}
\end{equation}
\begin{equation}
    t'^{-\nu_{\perp}} = b^{-2\nu_{\perp}}t^{-\nu_{\perp}} \sim \xi_{\perp}'(t') = b^{-1/2}\xi_{\perp}(t) \sim b^{-1/2}t^{-\nu_{\perp}} \rightarrow \nu_{\perp} = \frac{1}{4}
\end{equation}

\section{Upper critical dimension}
Let us consider the Landua-Ginzburg free energy:
\begin{equation}
    \mathcal{F}[m(x)] = \frac{1}{2}\left(\nabla m\right)^2 + tm^2 + um^4
\end{equation}
In class, we saw that for every natural $n > 2$, the prefactor $g_n$ of $m^n$ changes under RG flow to be:
\begin{equation}
    g_n' = g_n l^{-dn/2 + n + d}
\end{equation}
Applying this rule for $u = g_4$, we get:
\begin{equation}
    u(l) \sim l^{4 - d} 
\end{equation}
Let us take $d = 5$, since it's the lowest dimension above the UCD. This gives:
\begin{equation}
    u(l) \sim l^{-1}
\end{equation}
Therefore, for $l_0 = 100$, we get a $u(l_0) \ll 1$. We can now solve the ODE for $t$, which is the regualr gaussian term:
\begin{equation}
    \frac{dt}{dl} = 2lt
\end{equation}
Which gives:
\begin{equation}
    t = t_0 e^{l^2}
\end{equation}
And, plussing $l_0 = 100$ we get:
\begin{equation}
    t = e^{10000}t_0
\end{equation}
Which is much larger then the $u(l_0)$ term. This demonstrates that under RG flow for $d > 4$, $u$ grows vanihsingly small with respect to the other terms. Since we took the lowest dimension, this wil also be true for any dimension above the UCD $d > 4$.
\newline
The last term, in k-space, becomes:
\begin{equation}
    \int d^dr u m^4(r) = \frac{u}{(2\pi)^{4d}}\int d^dk_1 \cdots d^dk_4 \tilde{m}(k_1)\cdots\tilde{m}(k_4)(2\pi)^d\delta(k_1 + k_2 + k_3 + k_4)
\end{equation}
This, unlike the $m^2$ term, couples the $k = 0$ terms with the $k \ne 0$ terms. When $u$ is small, this coupling can be neglected, and we can write a theory for $k = 0$ alone. It will include the two original terms for $k = 0$, and the term in the above integral, with all four momenta being zero. Thsi gives:
\begin{equation}
    \mathcal{F}[\tilde{m}(0)] = \frac{1}{2}t\tilde{m}^2(0) + u\tilde{m}^4(0)
\end{equation}
This theory for $k = 0$ takes the same structure as a Landau theory, and so would give the same cirtical exponents as the Landau theory.
\newline
Now let us consider the follwing Landau-Ginzburg free energy:
\begin{equation}
    \mathcal{F}[m(x)] = \frac{1}{2}\left(\nabla m\right)^2 + tm^2 + vm^6
\end{equation}
and find its UCD using the different approaches.
\subsection{Ginzburg argument}
Let us calcualte the Ginzburg ratio:
\begin{equation}
    R = \frac{\int_{|\boldsymbol{r}| < \xi} d^dr\langle m(\boldsymbol{r})^2 \rangle}{\xi^d m_0^2}
\end{equation}
First, we start by calculating the nominator, which is the fluctuations of $m$. We go above the transition where $m = 0$, allowing us to ignore the $m^6$ term, and get the same gaussian approximation we got for the $m^4$ power in class. This gives the same nominator:
\begin{equation}
    \int_{|\boldsymbol{r}| < \xi} d^dr\langle m(\boldsymbol{r})^2 \rangle \sim t^{-1}
\end{equation}
The denominator, however, is different. Let us minimize the free energy as a function of $m$ to find $m_0$:
\begin{equation}
    \frac{dF}{dm} = 2mt + 6vm^5 \rightarrow m \sim t^{1/4}
\end{equation}
So, the denominator is:
\begin{equation}
    \xi^d m_0^2 \sim t^{\frac{1 - d}{2}}
\end{equation}
\begin{equation}
    R \sim \frac{t^{-1}}{t^{\frac{1 - d}{2}}} = t^{\frac{1}{2}(d - 6)}
\end{equation}
Which gives the UCD $d = 6$.
\subsection{K-space RG}
Let us now find the same from K-space RG to first order. The K-space free energy is:
\begin{equation}
    \mathcal{F}[\tilde{m}(k)] = \int^\Lambda \frac{d^dk}{(2\pi)^d}\frac{1}{2}(k^2 + t)\tilde{m}(k)\tilde{m}(-k) + \frac{v}{(2\pi)^{5d}}\int^\Lambda d^dk_1 \cdots d^dk_6 \tilde{m}(k_1) \cdots \tilde{m}(k_6) \delta(k_1 + \cdots + k_6)
\end{equation}
This gives the same partiaion function as the $m^4$ term we saw in class:
\begin{equation}
    Z = \int \mathcal{D}[\tilde{m}^{-}] \left[
        e^{-\frac{1}{2}\int_0^{\Lambda/b}\frac{d^dk}{(2\pi)^d}(t+k^2)\tilde{m}^{-}(k)\tilde{m}^{+}(-k)}
        \int \mathcal{D}[\tilde{m}^{+}]
        e^{-\frac{1}{2}\int_{\Lambda/b}^\Lambda\frac{d^dk}{(2\pi)^d}(t+k^2)\tilde{m}^{-}(k)\tilde{m}^{+}(-k)}
        e^{-F_I(\tilde{m}^{+}\tilde{m}^{-})}
    \right]
\end{equation}
Only with the interaction:
\begin{equation}
    F_I(\tilde{m}^{+}, \tilde{m}^{-}) = \frac{v}{(2\pi)^{5d}}\int_0^\Lambda d^dk_1 \cdots d^dk_6 \tilde{m}(k_1) \cdots \tilde{m}(k_6) \delta(k_1 + \cdots + k_6)
\end{equation}
The inner integral is again the same as doing a distribution
\begin{equation}
    \int \mathcal{D}[\tilde{m}^{+}]
        e^{-\frac{1}{2}\int_{\Lambda/b}^\Lambda\frac{d^dk}{(2\pi)^d}(t+k^2)\tilde{m}^{-}(k)\tilde{m}^{+}(-k)}
        e^{-F_I(\tilde{m}^{+}\tilde{m}^{-})}
         = Z_{m^{+}} \langle e^{-F_I} \rangle
\end{equation}
And to first order:
\begin{equation}
    \langle e^{-F_I} \rangle = 1 - \langle F_I \rangle
\end{equation}
So, evaluating $\langle F_I \rangle$, we get three surviving terms - the terms with two, four, and six $m^{+}$:
\begin{equation}
\begin{gathered}
    \tilde{m}_1^{-}\tilde{m}_2^{-}\tilde{m}_3^{-}
    \tilde{m}_4^{-}\langle\tilde{m}_5^{+}\tilde{m}_6^{+}\rangle
    \qquad
    \tilde{m}_1^{-}\tilde{m}_2^{-}\langle\tilde{m}_3^{+}
    \tilde{m}_4^{+}\tilde{m}_5^{+}\tilde{m}_6^{+}\rangle
    \qquad
    \langle\tilde{m}_1^{+}\tilde{m}_2^{+}\tilde{m}_3^{+}
    \tilde{m}_4^{+}\tilde{m}_5^{+}\tilde{m}_6^{+}\rangle
\end{gathered}
\end{equation}
The last term is just a constant, and so we focus on the first two terms. The first term gives the gaussian integral:
\begin{equation}
    \langle\tilde{m}_5^{+}\tilde{m}_6^{+}\rangle = \frac{(2\pi)^d \delta(k_5 + k_6)}{t + k_5^2}
\end{equation}
The delta function gives $k_5 = -k_6$, which together with the original delta function also gives $k_1 = -k_2$. Pluggin this into the integral over $\tilde{m}_1^{-}\tilde{m}_2^{-}$, we get the first term contribution:
\begin{equation}
    6v\left(\int \frac{d^dk}{(2\pi)^d} \tilde{m}^{-}(k)\tilde{m}^{-}(-k)\right)
    \left(
        \int \frac{d^dq}{(2\pi)^d}\frac{1}{t + q^2}
    \right)
\end{equation}
The second term, however, is a bit more complex, since it has an average over four $m^{+}$. This means we have to split it using Wick's theorem. We did it in class, and the result was the same as the two $m^{+}$, only with the denominator squared and the prefactor 36:
\begin{equation}
    36v\left(\int \frac{d^dk}{(2\pi)^d} \tilde{m}^{-}(k)\tilde{m}^{-}(-k)\right)
    \left(
        \int \frac{d^dq}{(2\pi)^d}\frac{1}{\left(t + q^2\right)^2}
    \right)
\end{equation}
And so, up to a constant:
\begin{equation}
    \langle F_I \rangle = \left(
        \int \frac{d^dk}{(2\pi)^d} \tilde{m}^{-}(k)\tilde{m}^{-}(-k)
    \right)
    \left(
        6 \int \frac{d^dq}{(2\pi)^d}\frac{1}{t + q^2} + 36 \int \frac{d^dq}{(2\pi)^d}\frac{1}{\left(t + q^2\right)^2}
    \right)
\end{equation}
And so, the flow of $t$ is now given by:
\begin{equation}
    \bar{t} = t + \int_{\Lambda/b}^\Lambda \frac{d^dq}{(2\pi)^d} \frac{ 12}{t + q^2} + \frac{72}{\left(t + q^2\right)^2}
\end{equation}
Let us perform the rescaling and renormalization steps:
\begin{equation}
    k' = bk
    \qquad
    \tilde{m}' = b^{\omega}\tilde{m}
\end{equation}
This gives the free energy:
\begin{equation}
\begin{gathered}
    \mathcal{F}[\tilde{m}(k')] = \frac{1}{2}\int_{0}^{\Lambda/b} \frac{d^dk'}{(2\pi)^d}b^{-d}(\bar{t} + b^{-2}k'^2)b^{-2\omega}\tilde{m}'(k)\tilde{m}'(-k) + \\ +
    v\int_{0}^{\Lambda/b} \frac{d^dk_1'\cdots d^dk_6'}{(2\pi)^{5d}}b^{-5d}b^{-6\omega}\tilde{m}'(k_1)\cdots\tilde{m}'(k_6)
\end{gathered}
\end{equation}
Again the $k^2$ coefficiant must be 1:
\begin{equation}
    b^{-2d-2-2\omega} = 1 \rightarrow \omega = -\frac{d}{2} - 1
\end{equation}
Which gives the recursion relations:
\begin{equation}
\begin{gathered}
    t' = b^2\left(t + \int_{\Lambda/b}^\Lambda \frac{d^dq}{(2\pi)^d} \frac{ 12}{t + q^2} + \frac{72}{\left(t + q^2\right)^2}\right)
    \\
    v' = b^{-5d - 6\omega}v = b^{6 - d}v
\end{gathered}
\end{equation}
This gives the UCD $d = 6$.
\end{document}